%! Orthogonal Matrices and Gram--Schmidt — Unit 4, Lesson 4.4 — Exit Ticket
\documentclass[10pt]{article}
\usepackage{linalg-article}
\usepackage{linalg-boxes}
\newcommand{\TallMath}[1]{$\displaystyle #1\rule[-1.4em]{0pt}{3.2em}$}
\newcommand{\vv}{\mathbf{v}}
\newcommand{\xx}{\mathbf{x}}
\newcommand{\bb}{\mathbf{b}}
\newcommand{\avec}{\mathbf{a}}
\newcommand{\qq}{\mathbf{q}}
\newcommand{\T}{\mathsf{T}}

\begin{document}

\pageheader{Unit 4, Lesson 4.4}{Exit Ticket}
\namedateperiod

\vspace{0.15in}

No notes. Show your reasoning.

\begin{enumerate}[leftmargin=1.8em, itemsep=16pt]
  \item \textbf{Orthonormal?} Are $q_1=\tfrac13(2,2,1)$ and $q_2=\tfrac13(2,-1,-2)$ orthonormal? Compute the
        three dot products $q_1\!\cdot\!q_1$, $q_2\!\cdot\!q_2$, $q_1\!\cdot\!q_2$ and decide.
        \writelines{2}
  \item \textbf{First Gram--Schmidt step.} Normalize $\avec=(1,1,1,1)$ to a unit vector $q_1=\avec/\|\avec\|$.
        \writelines{2}
  \item \textbf{Why it helps.} If $Q$ has orthonormal columns, what is $Q^{\T}Q$? Explain in one or two
        sentences why that makes the least-squares answer $\hat{\xx}=Q^{\T}\bb$ --- with no system to solve.
        \writelines{2}
\end{enumerate}

\end{document}
